% ============================================================
% 4_D_student.tex — 学生端功能需求 + 交互流程
% 编写人：D  日期：2026-07-15
% ============================================================

\section{学生端功能需求}

\subsection{学生端功能模块划分}

学生端是本系统面向论文撰写者（学生）的核心入口，承担论文创建、编辑、格式化、导出全流程操作。学生端功能划分为以下六大子模块：

\begin{enumerate}
    \item \textbf{账号与个人中心}：注册、登录、个人信息维护、密码修改；
    \item \textbf{论文项目管理}：新建论文、论文列表查看、论文信息编辑、论文删除；
    \item \textbf{章节在线编辑}：富文本章节编写、章节拖拽排序、章节增删、自动保存；
    \item \textbf{内容元素管理}：参考文献录入与格式化引用、图片上传与插入、表格创建与编辑；
    \item \textbf{模板选择与预览}：模板浏览、模板套用、论文实时预览；
    \item \textbf{导出下载}：论文导出为DOCX、PDF格式文件，下载至本地。
\end{enumerate}

\subsection{学生账号与个人中心}

\reqitem{FR-D01}{学生可自主注册账号，填写学号、姓名、邮箱、密码，系统校验学号唯一性及邮箱格式合法性，注册成功后自动跳转登录页。}
\reqitem{FR-D02}{学生使用学号+密码登录系统，后端返回JWT Token，前端存入localStorage，后续请求统一携带Token鉴权；Token过期后自动跳转登录页。}
\reqitem{FR-D03}{学生登录后进入个人中心，可查看和修改姓名、邮箱、手机号等个人信息；修改密码需先输入原密码验证。}
\reqitem{FR-D04}{学生可查看自己创建的全部论文列表，显示论文标题、创建时间、最后修改时间、当前字数统计。}

\subsection{论文项目管理}

\reqitem{FR-D05}{学生点击"新建论文"，进入论文创建向导：（1）输入论文标题；（2）选择所属学院；（3）选择论文模板；（4）确认创建，系统自动按模板生成空白论文骨架（含预设章节结构）。}
\reqitem{FR-D06}{学生在论文列表页可对每篇论文执行：打开编辑、重命名、查看预览、删除操作。删除论文需二次弹窗确认，确认后论文及所有章节数据永久删除。}
\reqitem{FR-D07}{论文列表支持按创建时间、最后修改时间排序，支持按论文标题关键词搜索过滤。}

\subsection{章节在线编辑}

\reqitem{FR-D08}{进入论文编辑器后，左侧显示章节树形导航，右侧为富文本编辑区域。点击章节名称切换编辑目标章节，编辑区同步刷新。}
\reqitem{FR-D09}{编辑器提供完整富文本工具栏：加粗、斜体、下划线、删除线、标题层级（H1-H4）、文字颜色/背景色、段落对齐、有序/无序列表、缩进、撤销/重做。}
\reqitem{FR-D10}{支持在章节正文中通过工具栏按钮或快捷键插入：参考文献引用标记、图片、表格、脚注。引用标记在导出时自动替换为格式化引用编号。}
\reqitem{FR-D11}{学生可在章节树中右键或拖拽调整章节顺序，支持添加同级章节、添加子章节、删除章节、重命名章节。删除章节时若包含子章节，需提示用户一并删除。}
\reqitem{FR-D12}{编辑器每30秒自动保存当前编辑内容至后端，同时监听页面关闭事件，若存在未保存修改则弹窗提示用户。学生也可手动点击"保存"按钮立即保存。}

\subsection{参考文献管理}

\reqitem{FR-D13}{学生可进入论文参考文献管理面板，手动添加参考文献条目，填写：作者、标题、期刊/出版社、年份、卷号/期号、页码、DOI/URL。}
\reqitem{FR-D14}{系统支持常见参考文献类型：期刊论文[J]、会议论文[C]、书籍[M]、学位论文[D]、网络资源[EB/OL]。学生选择类型后，表单字段动态调整。}
\reqitem{FR-D15}{参考文献列表支持编辑、删除、拖拽排序。在正文中通过插入引用标记并选择对应文献，系统按出现顺序自动编号。}
\reqitem{FR-D16}{导出DOCX/PDF时，系统自动在文末生成格式化参考文献列表，格式遵循GB/T 7714-2015国家标准。}

\subsection{模板选择与预览}

\reqitem{FR-D17}{新建论文时或编辑过程中，学生可浏览管理员发布的所有可用模板，按学院筛选。每个模板展示封面缩略图、适用学院、页眉页脚样式说明。}
\reqitem{FR-D18}{学生选择模板后，系统实时将当前论文内容套用新模板样式进行预览，包括字体、行距、页边距、页眉页脚。预览满意后学生确认切换模板。}
\reqitem{FR-D19}{论文编辑过程中，学生可随时点击"预览"按钮，在独立页面查看完整论文排版效果，预览页与实际导出版本视觉效果一致。}

\subsection{导出下载}

\reqitem{FR-D20}{学生点击"导出"，弹出导出选项对话框：选择导出格式（DOCX / PDF），选择导出范围（全部章节 / 指定章节范围），确认后后端生成文件。}
\reqitem{FR-D21}{导出任务提交后，前端显示生成进度（排队中→生成中→完成），完成后自动触发浏览器下载。若生成失败，显示具体错误原因并允许重试。}
\reqitem{FR-D22}{系统保留最近3次导出记录，学生可在导出历史中重新下载之前生成的文件，避免重复生成。}

% ============================================================
\section{学生端核心交互流程}

\subsection{注册登录流程}

\begin{enumerate}
    \item 学生访问系统首页，未登录状态下显示登录/注册入口及系统简介；
    \item 首次使用点击"注册"，填写学号、姓名、邮箱、密码、确认密码；
    \item 前端校验：学号非空且为数字、姓名非空、邮箱格式合法、密码长度≥6位、两次密码一致；
    \item 提交注册请求至 \texttt{POST /api/v1/auth/register}，后端校验学号唯一性，返回注册结果；
    \item 注册成功自动跳转登录页，输入学号+密码提交至 \texttt{POST /api/v1/auth/login}；
    \item 登录成功返回JWT Token，前端存储并跳转至论文列表首页；
    \item 登录失败提示具体错误（账号不存在 / 密码错误 / 账号已禁用）。
\end{enumerate}

\subsection{论文创建与编辑主流程}

\begin{enumerate}
    \item 学生在论文列表页点击"新建论文"，弹出创建向导对话框；
    \item 步骤一：输入论文标题（必填，限200字），选择所属学院（下拉列表）；
    \item 步骤二：系统展示可用模板列表，学生点击模板卡片预览样式；
    \item 步骤三：确认创建，后端创建论文记录，根据所选模板初始化章节结构；
    \item 前端跳转至论文编辑器页面（路由 \texttt{/editor/:paperId}），左侧加载章节树，右侧加载第一个章节内容；
    \item 学生在编辑器中进行富文本编辑，30秒自动保存或手动保存；
    \item 学生点击"预览"查看排版效果，新标签页展示渲染后的HTML；
    \item 学生点击"导出"，选择格式和范围，提交导出任务，等待生成完成下载文件。
\end{enumerate}

\subsection{参考文献插入流程}

\begin{enumerate}
    \item 学生在编辑器工具栏点击"插入引用"按钮，弹出参考文献管理面板；
    \item 面板分为两栏：左侧为已添加的参考文献列表，右侧为添加/编辑表单；
    \item 学生可新增、编辑、删除文献条目，操作实时保存至后端；
    \item 学生选中一条或多条文献，点击"插入引用标记"，系统在编辑器光标位置插入 [1] 占位标记；
    \item 正文编辑过程中引用标记与参考文献列表保持同步：删除文献时提示正文中存在引用，新增文献后引用编号自动重排；
    \item 导出时后端遍历正文引用标记，按出现顺序编号，文末生成标准格式参考文献列表。
\end{enumerate}

\subsection{论文导出流程}

\begin{enumerate}
    \item 学生点击编辑器顶部"导出"按钮，弹出导出配置对话框；
    \item 选择导出格式：DOCX / PDF（单选）；
    \item 选择导出范围：全部章节 / 自定义范围（多选章节）；
    \item 可选勾选：是否包含封面页、是否包含目录、是否包含参考文献列表；
    \item 提交导出请求，后端返回任务ID，前端轮询获取进度；
    \item 状态流转：PENDING → PROCESSING → COMPLETED / FAILED；
    \item 生成完成返回文件下载URL，前端自动触发下载；
    \item 生成失败显示错误阶段及原因，允许重新导出。
\end{enumerate}

% ============================================================
\section{学生端界面原型说明}

\begin{itemize}
    \item \textbf{登录/注册页}：居中卡片式布局，左侧系统Logo及简介，右侧表单区域。登录与注册通过Tab切换，避免页面跳转。
    \item \textbf{论文列表页（首页）}：顶部导航栏含Logo、搜索框、用户头像下拉菜单（个人中心、退出登录）。主体区域为论文卡片网格或列表视图，每张卡片显示论文标题、模板名称、最后修改时间、字数。右上角"新建论文"按钮。
    \item \textbf{论文编辑器页}：三栏布局——左侧可收起的章节树面板（支持拖拽排序），中间主编辑区（富文本编辑器+顶部工具栏），右侧可收起的内容面板（参考文献管理、模板预览切换）。底部状态栏显示自动保存状态和字数统计。
    \item \textbf{个人中心页}：Tab切换个人信息编辑、修改密码、导出历史三个子页面。
\end{itemize}
